% F2: o ciclo do agente nos braços de consulta e de ferramentas. Diagrama em TikZ.
\vspace{2mm}
\noindent\begin{minipage}{\columnwidth}
\centering\small
\begin{tikzpicture}[
  box/.style={draw, rounded corners=2pt, align=center, inner sep=3pt, minimum height=7mm},
  arr/.style={-{Stealth[length=2mm]}}]
\node[box, text width=0.24\columnwidth] (q) at (-0.36\columnwidth,0) {Pergunta};
\node[box, text width=0.24\columnwidth] (m) at (0,0) {Modelo};
\node[box, text width=0.24\columnwidth] (a) at (0.36\columnwidth,0) {Resposta};
\node[box, text width=0.5\columnwidth] (t) at (0,-1.8) {Ferramenta};
\node[box, text width=0.5\columnwidth] (g) at (0,-3.4) {Grafo Neo4j};
\draw[arr] (q) -- (m);
\draw[arr] (m) -- (a);
\draw[arr] ([xshift=-4mm]m.south) -- node[left] {chamada} ([xshift=-4mm]t.north);
\draw[arr] ([xshift=4mm]t.north) -- node[right] {resultado} ([xshift=4mm]m.south);
\draw[arr] ([xshift=-4mm]t.south) -- ([xshift=-4mm]g.north);
\draw[arr] ([xshift=4mm]g.north) -- ([xshift=4mm]t.south);
\node[align=left] at (0.385\columnwidth,-1.8) {até \nMaxCalls{}\\ chamadas};
\end{tikzpicture}
\captionof{figure}{Ciclo do agente em \armTextToCypher{} e \armGraphTools{}. No primeiro, a ferramenta executa no grafo a consulta Cypher que o modelo escreve; no segundo, o modelo escolhe uma das \nTools{} ferramentas e seus parâmetros, e ela executa uma consulta fixa.}\label{fig:agente}
\end{minipage}\par\smallskip
